\clearpage
\phantomsection
\chapter{CÀI ĐẶT VÀ TRIỂN KHAI}

\section{Môi trường phát triển}

Hệ thống được phát triển bằng ngôn ngữ C\# trên nền tảng .NET 9.0, đảm bảo hiệu năng cao và khả năng chạy đa nền tảng.
\begin{itemize}
    \item \textbf{Ngôn ngữ lập trình:} C\# 13.0
    \item \textbf{Framework:} .NET 9.0
    \item \textbf{IDE:} Visual Studio Code 1.107.0
    \item \textbf{Quản lý phiên bản:} Git và GitHub
    \item \textbf{Định dạng dữ liệu:} JSON được sử dụng để lưu trữ toàn bộ dữ liệu game (thẻ bài, kẻ địch, di vật) và cấu hình thuật toán.
\end{itemize}

\subsection{Cấu trúc dự án}
Dự án được tổ chức theo cấu trúc module rõ ràng để dễ dàng quản lý và mở rộng.
\begin{lstlisting}[language=bash, caption={Cấu trúc thư mục dự án}]
RoguelikeDeckBuilder/
|-- src/
|   |-- Roguelike/
|   |   |-- Core/           # Core game engine
|   |   |   |-- AI/
|   |   |   |-- Combat/
|   |   |   |-- Data/
|   |   |   |-- Map/
|   |   |-- Optimization/   # Optimization modules
|   |   |   |-- GASingleObjective/
|   |   |   |-- NSGA2MultiObjective/
|   |   |-- Program.cs      # Entry point
|   |   |-- GameData.json   # Game data
\end{lstlisting}

\section{Triển khai các module cốt lõi}

\subsection{Core Game Layer}
Tầng lõi của game được hiện thực hóa qua các lớp C\# chính:
\begin{itemize}
    \item \texttt{GameController.cs:} Lớp điều phối chính, xử lý các luồng trạng thái của game (bản đồ, chiến đấu, sự kiện, phần thưởng).
    \item \texttt{CombatManager.cs:} Quản lý logic chiến đấu theo lượt, bao gồm việc chơi bài, hành động của kẻ địch, áp dụng hiệu ứng và kiểm tra điều kiện thắng/thua.
    \item \texttt{DeckManager.cs:} Quản lý các chồng bài (rút, bỏ, trên tay) và các thao tác như xào bài, rút bài.
    \item \texttt{MapGenerator.cs:} Sử dụng thuật toán sinh ngẫu nhiên có ràng buộc để tạo ra bản đồ các phòng hợp lệ cho mỗi lượt chơi.
\end{itemize}

\subsection{AI Agent}
Lớp \texttt{HeuristicPlayerAI.cs} hiện thực hóa giao diện \texttt{IPlayerAgent}, cung cấp logic cho các quyết định của AI tại mỗi trạng thái của game. Các hàm tính điểm được sử dụng để đánh giá nước đi tốt nhất trong chiến đấu và trên bản đồ.

\begin{lstlisting}[language=C, caption={Logic ra quyết định trong chiến đấu của AI agent}, label={lst:ai_combat_decision}]
public CombatDecision GetCombatDecision(GameRun run)
{
    var hero = run.TheHero;
    var enemies = run.CurrentCombat.Enemies
        .Where(e => e.CurrentHealth > 0).ToList();
    
    // Tinh toan sat thuong sap nhan
    int totalIncomingDamage = CalculateIncomingDamage(run, enemies);
    int neededBlock = Math.Max(0, totalIncomingDamage - hero.Block);

    // UU TIEN 1: SINH TON - Choi the phong thu neu can
    if (neededBlock > 0)
    {
        var defensiveMove = FindBestDefensiveMove(
            run, hero, enemies, neededBlock);
        if (defensiveMove.HasValue)
            return defensiveMove.Value;
    }

    // UU TIEN 2: KET LIEU - Tim nuoc danh ket lieu ke dich
    var lethalMove = FindLethalMove(run, hero, enemies);
    if (lethalMove.HasValue)
        return lethalMove.Value;

    // UU TIEN 3: NUOC DI TOT NHAT - Choi the co diem cao nhat
    return GetBestGeneralMove(run, hero, enemies);
}
\end{lstlisting}

\section{Triển khai các module tối ưu hóa}
\subsection{GA đơn mục tiêu}
\begin{itemize}
    \item \texttt{GeneticAlgorithm.cs:} Lớp chính điều khiển vòng lặp tiến hóa, bao gồm chọn lọc, lai ghép và đột biến.
    \item \texttt{BalanceGenome.cs:} Biểu diễn bộ gen dưới dạng các `Dictionary` chứa các hệ số nhân cho từng tham số game.
    \item \texttt{FitnessEvaluator.cs:} Tính toán một giá trị fitness duy nhất từ kết quả của nhiều lượt mô phỏng.
\end{itemize}
\subsection{NSGA-II đa mục tiêu}
\begin{itemize}
    \item \texttt{ImprovedGeneticAlgorithm.cs:} Lớp điều khiển chính, sử dụng \texttt{NSGA2Optimizer}.
    \item \texttt{NSGA2Optimizer.cs:} Hiện thực hóa các thuật toán cốt lõi của NSGA-II:
        \begin{itemize}
            \item \texttt{FastNonDominatedSort}: Phân loại quần thể thành các Pareto front.
            \item \texttt{CalculateCrowdingDistance}: Tính toán khoảng cách đám đông để duy trì sự đa dạng.
            \item Sử dụng các toán tử lai ghép (Simulated Binary Crossover - SBX) và đột biến (Polynomial Mutation) phù hợp cho các giá trị thực.
        \end{itemize}
    \item \texttt{HierarchicalGenome.cs:} Biểu diễn bộ gen phân cấp với các thuộc tính cho từng nhóm tham số.
    \item \texttt{MultiObjectiveFitness.cs:} Tính toán 3 giá trị fitness riêng biệt (Balance, Engagement, Coherence).
\end{itemize}
\begin{lstlisting}[language=C, caption={Thuật toán Fast Non-Dominated Sort của NSGA-II}, label={lst:nsga2_sort}]
public List<List<Individual>> FastNonDominatedSort(
    List<Individual> population)
{
    var fronts = new List<List<Individual>>();
    var dominationCount = new Dictionary<Individual, int>();
    var dominatedSolutions = new Dictionary<Individual, List<Individual>>();
    
    // Tim cac moi quan he domination
    var firstFront = new List<Individual>();
    foreach (var p in population)
    {
        foreach (var q in population)
        {
            if (p == q) continue;
            if (p.Fitness.Dominates(q.Fitness))
                dominatedSolutions[p].Add(q);
            else if (q.Fitness.Dominates(p.Fitness))
                dominationCount[p]++;
        }
        
        // Neu khong bi dominated boi ai => Front 1
        if (dominationCount[p] == 0)
        {
            p.Fitness.Rank = 1;
            firstFront.Add(p);
        }
    }
    fronts.Add(firstFront);
    
    // Xay dung cac front tiep theo
    int currentRank = 1;
    while (fronts[currentRank - 1].Any())
    {
        var nextFront = new List<Individual>();
        foreach (var p in fronts[currentRank - 1])
        {
            foreach (var q in dominatedSolutions[p])
            {
                dominationCount[q]--;
                if (dominationCount[q] == 0)
                {
                    q.Fitness.Rank = currentRank + 1;
                    nextFront.Add(q);
                }
            }
        }
        if (nextFront.Any())
        {
            fronts.Add(nextFront);
            currentRank++;
        }
        else break;
    }
    return fronts;
}
\end{lstlisting}

\begin{lstlisting}[language=C, caption={Tính toán Crowding Distance để duy trì diversity}, label={lst:crowding_distance}]
public void CalculateCrowdingDistance(List<Individual> front)
{
    int n = front.Count;
    if (n == 0) return;
    
    // Khoi tao distance = 0
    foreach (var ind in front)
        ind.Fitness.CrowdingDistance = 0;
    
    // Voi moi objective (Balance, Engagement, Coherence)
    var objectives = new Func<MultiObjectiveFitness, float>[]
    {
        f => f.BalanceScore,
        f => f.EngagementScore,
        f => f.CoherenceScore
    };
    
    foreach (var objective in objectives)
    {
        // Sap xep theo objective nay
        var sorted = front.OrderBy(ind => objective(ind.Fitness))
            .ToList();
        
        // Cac giai phap o bien duoc gan distance vo cung
        sorted[0].Fitness.CrowdingDistance = float.MaxValue;
        sorted[n-1].Fitness.CrowdingDistance = float.MaxValue;
        
        float range = objective(sorted[n-1].Fitness) 
                    - objective(sorted[0].Fitness);
        if (range < 1e-6) continue;
        
        // Tinh crowding distance cho cac giai phap giua
        for (int i = 1; i < n - 1; i++)
        {
            float distance = (objective(sorted[i+1].Fitness) - 
                             objective(sorted[i-1].Fitness)) / range;
            sorted[i].Fitness.CrowdingDistance += distance;
        }
    }
}
\end{lstlisting}

\begin{lstlisting}[language=C, caption={Hàm kiểm tra Pareto dominance giữa hai giải pháp}, label={lst:dominates}]
public bool Dominates(MultiObjectiveFitness other)
{
    // Constraint-domination: giai phap kha thiluon tot hon khong kha thi
    if (this.IsFeasible && !other.IsFeasible) return true;
    if (!this.IsFeasible && other.IsFeasible) return false;
    
    // Ca hai khong kha thi: so sanh so luong rang buoc vi pham
    if (!this.IsFeasible && !other.IsFeasible)
        return this.ConstraintViolations < other.ConstraintViolations;
    
    // Ca hai kha thi: Pareto domination chuan
    bool atLeastOneBetter = false;
    
    // Phai >= voi tat ca cac objective
    if (BalanceScore < other.BalanceScore) return false;
    if (EngagementScore < other.EngagementScore) return false;
    if (CoherenceScore < other.CoherenceScore) return false;
    
    // Va phai tot hon strict o it nhat 1 objective
    if (BalanceScore > other.BalanceScore) atLeastOneBetter = true;
    if (EngagementScore > other.EngagementScore) atLeastOneBetter = true;
    if (CoherenceScore > other.CoherenceScore) atLeastOneBetter = true;
    
    return atLeastOneBetter;
}
\end{lstlisting}

\begin{lstlisting}[language=C, caption={Tính toán Balance Score từ các metrics gameplay}, label={lst:balance_score}]
private float CalculateBalanceScore(List<SimulationStats> results)
{
    // Tinh cac metrics thuc te
    float actualWinRate = (float)results.Count(r => r.IsVictory) / results.Count;
    
    var victories = results.Where(r => r.IsVictory).ToList();
    float actualVictoryHp = victories.Any() 
        ? victories.Average(r => (float)r.FinalHeroHP / r.FinalHeroMaxHP)
        : 0f;
    
    var deaths = results.Where(r => !r.IsVictory).ToList();
    float actualFloorOnDeath = deaths.Any()
        ? deaths.Average(r => r.FinalFloor)
        : 0f;
    
    // Cac gia tri muc tieu
    const float targetWinRate = 0.45f;
    const float targetVictoryHp = 0.30f;
    const float targetFloorOnDeath = 10f;
    
    // Tinh chenh lech chuan hoa
    float wrDeviation = Math.Abs(actualWinRate - targetWinRate) / targetWinRate;
    float hpDeviation = Math.Abs(actualVictoryHp - targetVictoryHp) / targetVictoryHp;
    float floorDeviation = Math.Abs(actualFloorOnDeath - targetFloorOnDeath) 
                          / targetFloorOnDeath;
    
    // Ket hop voi trong so
    float totalDeviation = 
        0.50f * wrDeviation +      // Win rate quan trong nhat
        0.30f * hpDeviation +      // Victory HP
        0.20f * floorDeviation;    // Floor on death
    
    // Chuyen thanh diem (1.0 = hoan hao, 0.0 = sai lech lon)
    return Math.Max(0f, 1.0f - totalDeviation);
}
\end{lstlisting}

\section{Thách thức và giải pháp}
\begin{itemize}
    \item \textbf{Tối ưu hóa hiệu năng:} Việc chạy hàng trăm nghìn lượt mô phỏng đòi hỏi hiệu năng cao. Giải pháp là sử dụng \texttt{Parallel.ForEach} của .NET để chạy các mô phỏng song song, tận dụng tối đa các lõi CPU.
    \item \textbf{Không gian tìm kiếm lớn:} Với phương pháp GA đơn mục tiêu, không gian ~500 chiều là một thách thức lớn. Giải pháp NSGA-II đa mục tiêu với \texttt{HierarchicalGenome} đã giảm không gian này xuống còn ~50 chiều, giúp thuật toán hội tụ nhanh và hiệu quả hơn.
    \item \textbf{Đảm bảo tính nhất quán:} Các thay đổi ngẫu nhiên của GA có thể tạo ra các tham số vô lý. Giải pháp là áp dụng các ràng buộc trong quá trình đột biến và lai ghép, đồng thời thiết kế \texttt{HierarchicalGenome} để các thay đổi mang tính tương đối thay vì tuyệt đối.
\end{itemize}